\scanpage{25}
En effet, soit $\psi$ une autre polarisation, donc
$\widetilde\psi=\widetilde\varphi\circ u$, où $u\in\operatorname{End}(A)\otimes_{\mathbb Z}\mathbb Q$, \deleted{\unsure{la}} et la condition
$\widetilde\psi'=-\widetilde\psi$ (\unsure{jointe à|joint à|pour :} $\widetilde\varphi'=-\widetilde\varphi$) donne
$\underset{\substack{\shortparallel\\-u'\widetilde\varphi}}{u'\widetilde\varphi'}=-\widetilde\varphi u$
i.e. $u'=\widetilde\varphi u\widetilde\varphi^{-1}$. On veut
prouver [posant $\delta(\widetilde\varphi)=\sqrt{\deg\widetilde\varphi}$, $\alpha_p(\varphi)=\Delta(\widetilde\varphi)$]
\[
\begin{aligned}
\Delta(\unsure{\widetilde\psi|\widetilde\varphi})/\delta(\unsure{\widetilde\varphi|\widetilde\psi})&=\Delta(\widetilde\varphi)/\delta(\widetilde\varphi)\\
\shortparallel\qquad&\\[-1ex]
\Delta(\widetilde\varphi)\circ\Delta(u)/\delta(\widetilde\varphi)\,\delta(u)&
\end{aligned}
\]
qui équivaut à : $\Delta(u)/\delta(u)=1\unsure{{}_{\ell}|{}}$. Or par
\unsure{Serre|dual.|lemme} on a $\deleted{\Delta}\Delta(u)\Delta(u^*)=\underset{\substack{\shortparallel\\\deg u}}{\delta(u)^2}$

d'autre part, comme $u^*=\widetilde\varphi u\widetilde\varphi\deleted{\unsure{u^*}}^{\inserted{-1}}$, on a
ici $\Delta(u^*)=\Delta(u)$, donc \deleted{$\unsure{\Delta(u)^2=}$}
\[
\Delta(u)^2=\delta(u)^2
\]
d'où $\Delta(u)=\pm\delta(u)$. Il faut encore examiner
le signe, il faut \deleted{\unsure{pour dire|prendre}} dans un $\operatorname{End}(A)\otimes_{\mathbb Z}\mathbb Q$,
muni de l'involution $u\leadsto u'=\widetilde\varphi u\widetilde\varphi^{-1}$, prouver que
\deleted{$u=u'$}\inserted{$u>0$ (i.e. \deleted{\unsure{vérifiant}} $u=\sum v.v'$)} implique
$\Delta(u)=\uuline{+}\delta(u)$ [N.B. $\delta(u)$ est bien
défini, car $\deg u$ est bien un carré parfait si $u=u'$]

\deleted{Or l'identité $\Delta(u)^2=\delta(u)$ est une identité dans}
% archive sheet number: 24
